%
% 6.006 problem set 3 solutions template
%
\documentclass[12pt,twoside]{article}

\newcommand{\name}{}

\usepackage{amssymb}
\usepackage{amsmath}
\usepackage{graphicx}
\usepackage{latexsym}
\usepackage{times,url}
\usepackage{cprotect}
\usepackage{listings}
\usepackage{graphicx}
\usepackage[table]{xcolor}
\usepackage[letterpaper]{geometry}
\usepackage{tikz-qtree}
\usepackage{enumerate}

\newcommand{\profs}{Instructors: Erik Demaine, Jason Ku, and Justin Solomon}
\newcommand{\subj}{6.006}
\newcommand{\ttt}[1]{{\tt\small #1}}

\definecolor{dkgreen}{rgb}{0,0.6,0}
\definecolor{dkblue}{rgb}{.2,.2,1}
\definecolor{gray}{rgb}{0.5,0.5,0.5}
\definecolor{mauve}{rgb}{0.58,0,0.82}

\lstset{
  language=Python,
  aboveskip=1pc,
  belowskip=1pc,
  basicstyle={\footnotesize\ttfamily},
  numbers=left,
  showstringspaces=false,
  numberstyle={\tiny\color{gray}\ttfamily},
  keywordstyle={\color{dkblue}\ttfamily},
  commentstyle={\color{dkgreen}\ttfamily},
  stringstyle={\color{mauve}\ttfamily},
}

% \lstset{
%   language=Python,
%   aboveskip=1pc,
%   belowskip=1pc,
%   basicstyle={\bf\color{white}\ttfamily},
%   numbers=left,
%   showstringspaces=false,
%   numberstyle={\bf\small\color{lightgray}\ttfamily},
%   keywordstyle={\bf\color{cyan}\ttfamily},
%   commentstyle={\bf\color{green}\ttfamily},
%   stringstyle={\bf\color{mauve}\ttfamily},
% }

\tikzset{
  % every node/.style={minimum width=2em,draw,circle},
  % level 1/.style={sibling distance=2cm},
  level distance=1cm,
  edge from parent/.style=
  {draw,edge from parent path={(\tikzparentnode) -- (\tikzchildnode)}},
}
\usetikzlibrary{shapes}

\newif\ifHideSolutions
\newcommand{\solution}[1]{\color{dkgreen}\textbf{Solution: }#1\color{black}}
\newcommand{\commonmistakes}[1]{{\color{dkblue}\textbf{Common Mistakes: }#1}}
\newcommand{\rubric}[1]{\color{dkgreen}{\bf Rubric:} #1\color{black}}

% \HideSolutionsfalse
% \ifHideSolutions
%   \renewcommand{\solution}[1]{}
%   \renewcommand{\rubric}[1]{}
% \fi

\newlength{\toppush}
\setlength{\toppush}{2\headheight}
\addtolength{\toppush}{\headsep}

\newcommand{\htitle}[2]{\noindent\vspace*{-\toppush}\newline\parbox{\textwidth}
{\textit{Introduction to Algorithms: 6.006}\hfill\name\newline
Massachusetts Institute of Technology \hfill #2\newline
\profs\hfill #1 \\[-3.5ex]\newline
\mbox{}\hrulefill\mbox{}}\vspace*{1ex}\mbox{}\newline
\begin{center}{\Large\bf #1}\end{center}}

\newcommand{\handout}[2]{\thispagestyle{empty}
 \markboth{#1}{#1}
 \pagestyle{myheadings}\htitle{#1}{#2}}

\newcommand{\lecture}[3]{\thispagestyle{empty}
 \markboth{Lecture #1: #2}{Lecture #1: #2}
 \pagestyle{myheadings}\htitle{Lecture #1: #2}{#3}}

\newcommand{\htitlewithouttitle}[2]{\noindent\vspace*{-\toppush}\newline\parbox{6.5in}
{\textit{Introduction to Algorithms}\hfill#2\newline
Massachusetts Institute of Technology \hfill 6.006\newline
\profs\hfill Handout #1\vspace*{-.5ex}\newline
\mbox{}\hrulefill\mbox{}}\vspace*{1ex}\mbox{}\newline}

\newcommand{\handoutwithouttitle}[2]{\thispagestyle{empty}
 \markboth{Handout \protect\ref{#1}}{Handout \protect\ref{#1}}
 \pagestyle{myheadings}\htitlewithouttitle{\protect\ref{#1}}{#2}}

\newcommand{\exam}[2]{% parameters: exam name, date
 \thispagestyle{empty}
 \markboth{\hspace{1cm}\subj\ #1\hspace{1in}Name\hrulefill\ \ }%
          {\subj\ #1\hspace{1in}Name\hrulefill\ \ }
 \pagestyle{myheadings}\examtitle{#1}{#2}
 \renewcommand{\theproblem}{Problem \arabic{problemnum}}
}
\newcommand{\examsolutions}[3]{% parameters: handout, exam name, date
 \thispagestyle{empty}
 \markboth{Handout \protect\ref{#1}: #2}{Handout \protect\ref{#1}: #2}
% \pagestyle{myheadings}\htitle{\protect\ref{#1}}{#2}{#3}
 \pagestyle{myheadings}\examsolutionstitle{\protect\ref{#1}} {#2}{#3}
 \renewcommand{\theproblem}{Problem \arabic{problemnum}}
}
\newcommand{\examsolutionstitle}[3]{\noindent\vspace*{-\toppush}\newline\parbox{6.5in}
{\textit{Introduction to Algorithms}\hfill#3\newline
Massachusetts Institute of Technology \hfill 6.006\newline
%Singapore-MIT Alliance \hfill SMA5503\newline
\profs\hfill Handout #1\vspace*{-.5ex}\newline
\mbox{}\hrulefill\mbox{}}\vspace*{1ex}\mbox{}\newline
\begin{center}{\Large\bf #2}\end{center}}

\newcommand{\takehomeexam}[2]{% parameters: exam name, date
 \thispagestyle{empty}
 \markboth{\subj\ #1\hfill}{\subj\ #1\hfill}
 \pagestyle{myheadings}\examtitle{#1}{#2}
 \renewcommand{\theproblem}{Problem \arabic{problemnum}}
}

\makeatletter
\newcommand{\exambooklet}[2]{% parameters: exam name, date
 \thispagestyle{empty}
 \markboth{\subj\ #1}{\subj\ #1}
 \pagestyle{myheadings}\examtitle{#1}{#2}
 \renewcommand{\theproblem}{Problem \arabic{problemnum}}
 \renewcommand{\problem}{\newpage
 \item \let\@currentlabel=\theproblem
 \markboth{\subj\ #1, \theproblem}{\subj\ #1, \theproblem}}
}
\makeatother


\newcommand{\examtitle}[2]{\noindent\vspace*{-\toppush}\newline\parbox{6.5in}
{\textit{Introduction to Algorithms}\hfill#2\newline
Massachusetts Institute of Technology \hfill 6.006 Spring 2020\newline
%Singapore-MIT Alliance \hfill SMA5503\newline
\profs\hfill #1\vspace*{-.5ex}\newline
\mbox{}\hrulefill\mbox{}}\vspace*{1ex}\mbox{}\newline
\begin{center}{\Large\bf #1}\end{center}}

\newcommand{\grader}[1]{\hspace{1cm}\textsf{\textbf{#1}}\hspace{1cm}}

\newcommand{\points}[1]{[#1 points]\ }
\newcommand{\parts}[1]
{
  \ifnum#1=1
  (1 part)
  \else
  (#1 parts)
  \fi
  \ 
}

\newcommand{\bparts}{\begin{problemparts}}
\newcommand{\eparts}{\end{problemparts}}
\newcommand{\ppart}{\problempart}

%\newcommand{\lg} {lg\ }

\setlength{\oddsidemargin}{0pt}
\setlength{\evensidemargin}{0pt}
\setlength{\textwidth}{6.5in}
\setlength{\topmargin}{0in}
\setlength{\textheight}{8.5in}


\newcommand{\Spawn}{{\bf spawn} }
\newcommand{\Sync}{{\bf sync}}

\renewcommand{\cases}[1]{\left\{ \begin{array}{ll}#1\end{array}\right.}
\newcommand{\cif}[1]{\mbox{if $#1$}}
\newcommand{\cwhen}[1]{\mbox{when $#1$}}

\newcounter{problemnum}
\newcommand{\theproblem}{Problem \theproblemsetnum-\arabic{problemnum}}
\newenvironment{problems}{
        \begin{list}{{\bf \theproblem. \hspace*{0.5em}}}
        {\setlength{\leftmargin}{0em}
         \setlength{\rightmargin}{0em}
         \setlength{\labelwidth}{0em}
         \setlength{\labelsep}{0em}
         \usecounter{problemnum}}}{\end{list}}
\makeatletter
\newcommand{\problem}[1][{}]{\item \let\@currentlabel=\theproblem \textbf{#1}}
\makeatother

\newcounter{problempartnum}[problemnum]
\newenvironment{problemparts}{
        \begin{list}{{\bf (\alph{problempartnum})}}
        {\setlength{\leftmargin}{2.5em}
         \setlength{\rightmargin}{2.5em}
         \setlength{\labelsep}{0.5em}}}{\end{list}}
\newcommand{\problempart}{\addtocounter{problempartnum}{1}\item}

\newenvironment{truefalseproblemparts}{
        \begin{list}{{\bf (\alph{problempartnum})\ \ \ T\ \ F\hfil}}
        {\setlength{\leftmargin}{4.5em}
         \setlength{\rightmargin}{2.5em}
         \setlength{\labelsep}{0.5em}
         \setlength{\labelwidth}{4.5em}}}{\end{list}}

\newcounter{exercisenum}
\newcommand{\theexercise}{Exercise \theproblemsetnum-\arabic{exercisenum}}
\newenvironment{exercises}{
        \begin{list}{{\bf \theexercise. \hspace*{0.5em}}}
        {\setlength{\leftmargin}{0em}
         \setlength{\rightmargin}{0em}
         \setlength{\labelwidth}{0em}
         \setlength{\labelsep}{0em}
        \usecounter{exercisenum}}}{\end{list}}
\makeatletter
\newcommand{\exercise}{\item \let\@currentlabel=\theexercise}
\makeatother

\newcounter{exercisepartnum}[exercisenum]
%\newcommand{\problem}[1]{\medskip\mbox{}\newline\noindent{\bf Problem #1.}\hspace*{1em}}
%\newcommand{\exercise}[1]{\medskip\mbox{}\newline\noindent{\bf Exercise #1.}\hspace*{1em}}

\newenvironment{exerciseparts}{
        \begin{list}{{\bf (\alph{exercisepartnum})}}
        {\setlength{\leftmargin}{2.5em}
         \setlength{\rightmargin}{2.5em}
         \setlength{\labelsep}{0.5em}}}{\end{list}}
\newcommand{\exercisepart}{\addtocounter{exercisepartnum}{1}\item}


% Macros to make captions print with small type and 'Figure xx' in bold.
\makeatletter
\def\fnum@figure{{\bf Figure \thefigure}}
\def\fnum@table{{\bf Table \thetable}}
\let\@mycaption\caption
%\long\def\@mycaption#1[#2]#3{\addcontentsline{\csname
%  ext@#1\endcsname}{#1}{\protect\numberline{\csname 
%  the#1\endcsname}{\ignorespaces #2}}\par
%  \begingroup
%    \@parboxrestore
%    \small
%    \@makecaption{\csname fnum@#1\endcsname}{\ignorespaces #3}\par
%  \endgroup}
%\def\mycaption{\refstepcounter\@captype \@dblarg{\@mycaption\@captype}}
%\makeatother
\let\mycaption\caption
%\newcommand{\figcaption}[1]{\mycaption[]{#1}}

\newcounter{totalcaptions}
\newcounter{totalart}

\newcommand{\figcaption}[1]{\addtocounter{totalcaptions}{1}\caption[]{#1}}

% \psfigures determines what to do for figures:
%       0 means just leave vertical space
%       1 means put a vertical rule and the figure name
%       2 means insert the PostScript version of the figure
%       3 means put the figure name flush left or right
\newcommand{\psfigures}{0}
\newcommand{\spacefigures}{\renewcommand{\psfigures}{0}}
\newcommand{\rulefigures}{\renewcommand{\psfigures}{1}}
\newcommand{\macfigures}{\renewcommand{\psfigures}{2}}
\newcommand{\namefigures}{\renewcommand{\psfigures}{3}}

\newcommand{\figpart}[1]{{\bf (#1)}\nolinebreak[2]\relax}
\newcommand{\figparts}[2]{{\bf (#1)--(#2)}\nolinebreak[2]\relax}


\macfigures     % STATE

% When calling \figspace, make sure to leave a blank line afterward!!
% \widefigspace is for figures that are more than 28pc wide.
\newlength{\halffigspace} \newlength{\wholefigspace}
\newlength{\figruleheight} \newlength{\figgap}
\newcommand{\setfiglengths}{\ifnum\psfigures=1\setlength{\figruleheight}{\hruleheight}\setlength{\figgap}{1em}\else\setlength{\figruleheight}{0pt}\setlength{\figgap}{0em}\fi}
\newcommand{\figspace}[2]{\ifnum\psfigures=0\leavefigspace{#1}\else%
\setfiglengths%
\setlength{\wholefigspace}{#1}\setlength{\halffigspace}{.5\wholefigspace}%
\rule[-\halffigspace]{\figruleheight}{\wholefigspace}\hspace{\figgap}#2\fi}
\newlength{\widefigspacewidth}
% Make \widefigspace put the figure flush right on the text page.
\newcommand{\widefigspace}[2]{
\ifnum\psfigures=0\leavefigspace{#1}\else%
\setfiglengths%
\setlength{\widefigspacewidth}{28pc}%
\addtolength{\widefigspacewidth}{-\figruleheight}%
\setlength{\wholefigspace}{#1}\setlength{\halffigspace}{.5\wholefigspace}%
\makebox[\widefigspacewidth][r]{#2\hspace{\figgap}}\rule[-\halffigspace]{\figruleheight}{\wholefigspace}\fi}
\newcommand{\leavefigspace}[1]{\setlength{\wholefigspace}{#1}\setlength{\halffigspace}{.5\wholefigspace}\rule[-\halffigspace]{0em}{\wholefigspace}}

% Commands for including figures with macpsfig.
% To use these commands, documentstyle ``macpsfig'' must be specified.
\newlength{\macfigfill}
\makeatother
\newlength{\bbx}
\newlength{\bby}
\newcommand{\macfigure}[5]{\addtocounter{totalart}{1}
\ifnum\psfigures=2%
\setlength{\bbx}{#2}\addtolength{\bbx}{#4}%
\setlength{\bby}{#3}\addtolength{\bby}{#5}%
\begin{flushleft}
\ifdim#4>28pc\setlength{\macfigfill}{#4}\addtolength{\macfigfill}{-28pc}\hspace*{-\macfigfill}\fi%
\mbox{\psfig{figure=./#1.ps,%
bbllx=#2,bblly=#3,bburx=\bbx,bbury=\bby}}
\end{flushleft}%
\else\ifdim#4>28pc\widefigspace{#5}{#1}\else\figspace{#5}{#1}\fi\fi}
\makeatletter

\newlength{\savearraycolsep}
\newcommand{\narrowarray}[1]{\setlength{\savearraycolsep}{\arraycolsep}\setlength{\arraycolsep}{#1\arraycolsep}}
\newcommand{\normalarray}{\setlength{\arraycolsep}{\savearraycolsep}}

\newcommand{\hint}{{\bf Hint:\ }}

% Macros from /th/u/clr/mac.tex

\newcommand{\set}[1]{\left\{ #1 \right\}}
\newcommand{\abs}[1]{\left| #1\right|}
\newcommand{\card}[1]{\left| #1\right|}
\newcommand{\floor}[1]{\left\lfloor #1 \right\rfloor}
\newcommand{\ceil}[1]{\left\lceil #1 \right\rceil}
\newcommand{\ang}[1]{\ifmmode{\left\langle #1 \right\rangle}
   \else{$\left\langle${#1}$\right\rangle$}\fi}
        % the \if allows use outside mathmode,
        % but will swallow following space there!
\newcommand{\paren}[1]{\left( #1 \right)}
\newcommand{\bracket}[1]{\left[ #1 \right]}
\newcommand{\prob}[1]{\Pr\left\{ #1 \right\}}
\newcommand{\Var}{\mathop{\rm Var}\nolimits}
\newcommand{\expect}[1]{{\rm E}\left[ #1 \right]}
\newcommand{\expectsq}[1]{{\rm E}^2\left[ #1 \right]}
\newcommand{\variance}[1]{{\rm Var}\left[ #1 \right]}
\renewcommand{\choose}[2]{{{#1}\atopwithdelims(){#2}}}
\def\pmod#1{\allowbreak\mkern12mu({\rm mod}\,\,#1)}
\newcommand{\matx}[2]{\left(\begin{array}{*{#1}{c}}#2\end{array}\right)}
\newcommand{\Adj}{\mathop{\rm Adj}\nolimits}

\newtheorem{theorem}{Theorem}
\newtheorem{lemma}[theorem]{Lemma}
\newtheorem{corollary}[theorem]{Corollary}
\newtheorem{xample}{Example}
\newtheorem{definition}{Definition}
\newenvironment{example}{\begin{xample}\rm}{\end{xample}}
\newcommand{\proof}{\noindent{\em Proof.}\hspace{1em}}
\def\squarebox#1{\hbox to #1{\hfill\vbox to #1{\vfill}}}
\newcommand{\qedbox}{\vbox{\hrule\hbox{\vrule\squarebox{.667em}\vrule}\hrule}}
\newcommand{\qed}{\nopagebreak\mbox{}\hfill\qedbox\smallskip}
\newcommand{\eqnref}[1]{(\protect\ref{#1})}

%%\newcommand{\twodots}{\mathinner{\ldotp\ldotp}}
\newcommand{\transpose}{^{\mbox{\scriptsize \sf T}}}
\newcommand{\amortized}[1]{\widehat{#1}}

\newcommand{\punt}[1]{}

%%% command for putting definitions into boldface
% New style for defined terms, as of 2/23/88, redefined by THC.
\newcommand{\defn}[1]{{\boldmath\textit{\textbf{#1}}}}
\newcommand{\defi}[1]{{\textit{\textbf{#1\/}}}}

\newcommand{\red}{\leq_{\rm P}}
\newcommand{\lang}[1]{%
\ifmmode\mathord{\mathcode`-="702D\rm#1\mathcode`\-="2200}\else{\rm#1}\fi}

%\newcommand{\ckt}[1]{\ifmmode\mathord{\mathcode`-="702D\sc #1\mathcode`\-="2200}\else$\mathord{\mathcode`-="702D\sc #1\mathcode`\-="2200}$\fi}
\newcommand{\ckt}[1]{\ifmmode \sc #1\else$\sc #1$\fi}

%% Margin notes - use \notesfalse to turn off notes.
\setlength{\marginparwidth}{0.6in}
\reversemarginpar
\newif\ifnotes
\notestrue
\newcommand{\longnote}[1]{
  \ifnotes
    {\medskip\noindent Note: \marginpar[\hfill$\Longrightarrow$]
      {$\Longleftarrow$}{#1}\medskip}
  \fi}
\newcommand{\note}[1]{
  \ifnotes
    {\marginpar{\tiny \raggedright{#1}}}
  \fi}


\newcommand{\reals}{\mathbbm{R}}
\newcommand{\integers}{\mathbbm{Z}}
\newcommand{\naturals}{\mathbbm{N}}
\newcommand{\rationals}{\mathbbm{Q}}
\newcommand{\complex}{\mathbbm{C}}

\newcommand{\oldreals}{{\bf R}}
\newcommand{\oldintegers}{{\bf Z}}
\newcommand{\oldnaturals}{{\bf N}}
\newcommand{\oldrationals}{{\bf Q}}
\newcommand{\oldcomplex}{{\bf C}}

\newcommand{\w}{\omega}                 %% for fft chapter

\newenvironment{closeitemize}{\begin{list}
{$\bullet$}
{\setlength{\itemsep}{-0.2\baselineskip}
\setlength{\topsep}{0.2\baselineskip}
\setlength{\parskip}{0pt}}}
{\end{list}}

% These are necessary within a {problems} environment in order to restore
% the default separation between bullets and items.
\newenvironment{normalitemize}{\setlength{\labelsep}{0.5em}\begin{itemize}}
                              {\end{itemize}}
\newenvironment{normalenumerate}{\setlength{\labelsep}{0.5em}\begin{enumerate}}
                                {\end{enumerate}}

%\def\eqref#1{Equation~(\ref{eq:#1})}
%\newcommand{\eqref}[1]{Equation (\ref{eq:#1})}
\newcommand{\eqreftwo}[2]{Equations (\ref{eq:#1}) and~(\ref{eq:#2})}
\newcommand{\ineqref}[1]{Inequality~(\ref{ineq:#1})}
\newcommand{\ineqreftwo}[2]{Inequalities (\ref{ineq:#1}) and~(\ref{ineq:#2})}

\newcommand{\figref}[1]{Figure~\ref{fig:#1}}
\newcommand{\figreftwo}[2]{Figures \ref{fig:#1} and~\ref{fig:#2}}

\newcommand{\liref}[1]{line~\ref{li:#1}}
\newcommand{\Liref}[1]{Line~\ref{li:#1}}
\newcommand{\lirefs}[2]{lines \ref{li:#1}--\ref{li:#2}}
\newcommand{\Lirefs}[2]{Lines \ref{li:#1}--\ref{li:#2}}
\newcommand{\lireftwo}[2]{lines \ref{li:#1} and~\ref{li:#2}}
\newcommand{\lirefthree}[3]{lines \ref{li:#1}, \ref{li:#2}, and~\ref{li:#3}}

\newcommand{\lemlabel}[1]{\label{lem:#1}}
\newcommand{\lemref}[1]{Lemma~\ref{lem:#1}} 

\newcommand{\exref}[1]{Exercise~\ref{ex:#1}}

\newcommand{\handref}[1]{Handout~\ref{#1}}

\newcommand{\defref}[1]{Definition~\ref{def:#1}}

% (1997.8.16: Victor Luchangco)
% Modified \hlabel to only get date and to use handouts counter for number.
%   New \handout and \handoutwithouttitle commands in newmac.tex use this.
%   The date is referenced by <label>-date.
%   (Retained old definition as \hlabelold.)
%   Defined \hforcelabel to use an argument instead of the handouts counter.

\newcounter{handouts}
\setcounter{handouts}{0}

\newcommand{\hlabel}[2]{%
\stepcounter{handouts}
{\edef\next{\write\@auxout{\string\newlabel{#1}{{\arabic{handouts}}{0}}}}\next}
\write\@auxout{\string\newlabel{#1-date}{{#2}{0}}}
}

\newcommand{\hforcelabel}[3]{%          Does not step handouts counter.
\write\@auxout{\string\newlabel{#1}{{#2}{0}}}
\write\@auxout{\string\newlabel{#1-date}{{#3}{0}}}}


% less ugly underscore
% --juang, 2008 oct 05
\renewcommand{\_}{\vrule height 0 pt depth 0.4 pt width 0.5 em \,}

% multiline framed box (will always extend to the far right edge; for a short single line, use \fbox directly)
% --zabel, fall 2018
\newcommand\framepar[1]{\fbox{\begin{minipage}{\linewidth}#1\end{minipage}}}


\newcommand{\theproblemsetnum}{3}

\title{6.006 Problem Set \theproblemsetnum}

\begin{document}

\handout{Problem Set \theproblemsetnum}

\setlength{\parindent}{0pt}
\medskip\hrulefill\medskip

{\bf Name:} Your Name

\medskip

{\bf Collaborators:} Name1, Name2

\medskip\hrulefill

%%%%%%%%%%%%%%%%%%%%%%%%%%%%%%%%%%%%%%%%%%%%%%%%%%%%%
% See below for common and useful latex constructs. %
%%%%%%%%%%%%%%%%%%%%%%%%%%%%%%%%%%%%%%%%%%%%%%%%%%%%%

% Some useful commands:
%$f(x) = \Theta(x)$
%$T(x, y) \leq \log(x) + 2^y + \binom{2n}{n}$
% {\tt code\_function}


% You can create unnumbered lists as follows:
%\begin{itemize}
%    \item First item in a list
%        \begin{itemize}
%            \item First item in a list
%                \begin{itemize}
%                    \item First item in a list
%                    \item Second item in a list
%                \end{itemize}
%            \item Second item in a list
%        \end{itemize}
%    \item Second item in a list
%\end{itemize}

% You can create numbered lists as follows:
%\begin{enumerate}
%    \item First item in a list
%    \item Second item in a list
%    \item Third item in a list
%\end{enumerate}

% You can write aligned equations as follows:
%\begin{align}
%    \begin{split}
%        (x+y)^3 &= (x+y)^2(x+y) \\
%                &= (x^2+2xy+y^2)(x+y) \\
%                &= (x^3+2x^2y+xy^2) + (x^2y+2xy^2+y^3) \\
%                &= x^3+3x^2y+3xy^2+y^3
%    \end{split}
%\end{align}

% You can create grids/matrices as follows:
%\begin{align}
%    A =
%    \begin{bmatrix}
%        A_{11} & A_{21} \\
%        A_{21} & A_{22}
%    \end{bmatrix}
%\end{align}

% You can include images and PDFs as follows:
% \includegraphics[width=0.5\textwidth]{img.jpg}

\begin{problems}

\problem  % Problem 1

\begin{problemparts}
\problempart % Problem 1a
Compute $h(k)=(10k+4)\bmod 7$ for each key, in order:
$47\to 5$, $61\to 5$, $36\to 0$, $52\to 6$, $56\to 4$, $33\to 5$, $92\to 0$.
Later keys are appended to the end of a slot's list, so the table is:

\begin{center}
\begin{tabular}{c|l}
slot & keys \\ \hline
0 & $36 \to 92$ \\
1 & \\
2 & \\
3 & \\
4 & $56$ \\
5 & $47 \to 61 \to 33$ \\
6 & $52$
\end{tabular}
\end{center}

\problempart % Problem 1b
Seven keys must go into seven slots with no collision, so $c\ge 7$ is needed (for $c<7$ the value
$(10k+4)\bmod c$ takes fewer than 7 values, so by pigeonhole two keys collide). Checking
$c=7,8,\dots$ in order, with the hash values of $(47,61,36,52,56,33,92)$:

\begin{center}
\begin{tabular}{c|c|c}
$c$ & hash values & distinct slots used \\ \hline
7 & $5,5,0,6,4,5,0$ & 4 \\
8 & $2,6,4,4,4,6,4$ & 3 \\
9 & $6,2,4,2,6,1,6$ & 4 \\
10 & $4,4,4,4,4,4,4$ & 1 \\
11 & $1,2,1,0,3,4,0$ & 5 \\
12 & $6,2,4,1,0,3,0$ & 6 \\
13 & $6,3,0,4,5,2,1$ & 7
\end{tabular}
\end{center}

Every $c$ from 7 to 12 has a collision, and for $c=13$ all seven hash values are different.
So the smallest value is $\mathbf{c=13}$.
\end{problemparts}

\newpage

\problem  % Problem 2

\begin{problemparts}
\problempart % Problem 2a
\textbf{They can guarantee it.} Choose two IDs with the same remainder mod $n$, for example
$k_1=1$ and $k_2=n+1$ (both are less than $u$ since $u\gg 2n$). For every $a,b$:
$ak_2+b = ak_1 + an + b \equiv ak_1+b \pmod n$, so $h_{ab}(k_1)=h_{ab}(k_2)$ for \emph{every}
function in the family. They are roommates with probability 1.

\problempart % Problem 2b
\textbf{They can guarantee it.} Here $h_a(k)=\left(\left\lfloor \frac{kn}{u}\right\rfloor + a\right)\bmod n$.
Choose $k_1=1$ and $k_2=2$. Since $u\gg 2n$, we have $0\le n/u < 2n/u < 1$, so
$\lfloor k_1 n/u\rfloor=\lfloor k_2 n/u\rfloor=0$. Both IDs then get the same value $a \bmod n$ for
every $a$, so $h_a(k_1)=h_a(k_2)$ for every function in the family. They are roommates with probability 1.

\problempart % Problem 2c
\textbf{They cannot guarantee it.} This is the universal hash family from lecture. For any two
different keys $k_1\ne k_2$, a random $h$ from a universal family satisfies
$\Pr[h(k_1)=h(k_2)]\le 1/m$, where $m$ is the number of slots. Here $m=n$, so Rony and Tiri
collide with probability at most $1/n$, whatever IDs they choose. Since $1/n<1$ (for $n\ge 2$), they
cannot guarantee being roommates; the best they can hope for is the universal bound $1/n$.
\end{problemparts}

\newpage

\problem  % Problem 3

\begin{problemparts}
\problempart % Problem 3a
\textbf{Radix sort, $\Theta(n)$.} The identifier has $16\lceil\log_4\sqrt n\rceil$ characters, and
$\log_4\sqrt n=\tfrac14\lg n$, so it has $O(\log n)$ characters of 8 bits each. Read the whole string as one
integer (comparing two fixed-length strings alphabetically is the same as comparing their integers).
That integer is less than
$2^{8\cdot 16\lceil\frac14\lg n\rceil}\le 2^{32\lg n+128}=2^{128}\,n^{32}$, a polynomial in $n$.
In the word-RAM a word has $w\ge\lg n$ bits, so the integer takes a constant number of words.
Radix sort in base $n$ handles integers up to $n^{c}$ in $c$ passes of counting sort, here
$O(n+n\log_n n^{32})=O(n)$ time.

\problempart % Problem 3b
\textbf{Counting sort, $\Theta(n)$.} Every age is an integer in $[0,800{,}000]$, a range of constant
size. Counting sort takes $\Theta(n+800{,}000)=\Theta(n)$ time. (Radix sort would also work.)

\problempart % Problem 3c
\textbf{Radix sort, $\Theta(n)$.} Multiply every thickness $m/n^3$ by $n^3$ to get the integer $m$,
with $0\le m\le 4n^3$. These integers are bounded by a polynomial in $n$, so radix sort in base
$n$ uses a constant number of counting-sort passes (about 4), for $\Theta(n)$ time. Sorting the
integers sorts the thicknesses in the same order.

\problempart % Problem 3d
\textbf{Merge sort, $\Theta(n\log n)$.} The only thing we can do with two slices is compare them
(in $O(1)$ time), and the number of comparisons needed to sort $n$ items in the comparison model is
$\Omega(n\log n)$. Merge sort does $O(n\log n)$ comparisons, so it is asymptotically optimal here;
no faster algorithm exists for this model.
\end{problemparts}

\newpage

\problem  % Problem 4

\begin{problemparts}
\problempart % Problem 4a
\textbf{Idea:} Each $b_i$ can only pair with one value, $r-b_i$. Use a hash table to look that
value up quickly instead of scanning all boxes.

\textbf{Algorithm.}
\begin{enumerate}
\item Insert every pair $(b_i, i)$ into a hash table $H$ keyed by $b_i$, with value $i$.
\item For each $i$, look up $r-b_i$ in $H$. If it is found with index $j$, and $j\ne i$ and
$|i-j|<n/10$, return True.
\item If no $i$ succeeds, return False.
\end{enumerate}

\textbf{Correctness.} The values are all different, so $r-b_i$ appears at most once in $B$, and the
table returns exactly the one box that could pair with box $i$. We check $j\ne i$ because $b_i=r/2$
would otherwise match itself. The algorithm tests every box, so it finds a close pair if one exists.

\textbf{Running time.} Building $H$ takes expected $O(n)$ time ($n$ inserts of expected $O(1)$ each),
and $n$ lookups take expected $O(1)$ each, so the total is expected $O(n)$.

\problempart % Problem 4b
\textbf{Idea:} With $r<n^2$, only values below $r$ matter, so we can radix sort them in $O(n)$ and
replace the hash table by a two-finger sweep on the sorted array.

\textbf{Algorithm.}
\begin{enumerate}
\item Make pairs $(b_i,i)$ and drop every pair with $b_i\ge r$ (the other number in the pair
must be positive, so such a box can never be part of a pair that sums to $r$). Call the remaining
array $B'$.
\item Every remaining value is less than $r<n^2$, so radix sort $B'$ by $b_i$ in base $n$
(two passes of counting sort) in $O(n)$ time. Call the sorted array $A$.
\item Set $i=0$ and $j=|A|-1$. While $i<j$:
\begin{itemize}
\item if $A[i].b+A[j].b=r$: if $|A[i].x-A[j].x|<n/10$ return True; otherwise set $i\gets i+1$ and
$j\gets j-1$ (neither box has another possible partner, because the values are distinct);
\item if the sum is less than $r$: $i\gets i+1$;
\item if the sum is greater than $r$: $j\gets j-1$.
\end{itemize}
\item Return False.
\end{enumerate}
(Here $A[i].b$ is the number of reams and $A[i].x$ the original position.)

\textbf{Correctness.} If $A[i].b+A[j].b<r$, then $A[i]$ plus any element at or before $A[j]$ is also
less than $r$, so $A[i]$ has no partner left and we can discard it. The case ``greater than $r$'' is
symmetric for $A[j]$. When the sum equals $r$, those two boxes are each other's only possible partner,
so after checking closeness both can be discarded. Every pair that sums to $r$ is therefore examined.

\textbf{Running time.} Filtering is $O(n)$ and radix sort is $O(n)$. Each loop step takes $O(1)$ time
and shrinks $j-i$ by at least one, so the loop runs at most $n$ times. Total worst-case $O(n)$.
\end{problemparts}

\newpage

\problem  % Problem 5

\begin{problemparts}
\problempart % Problem 5a
\textbf{Idea:} Two strings are anagrams exactly when they have the same letter-frequency table
(26 counts, one per letter). Store the frequency table of every length-$k$ substring of $A$ in
a hash table, with the number of substrings that have that table.

\textbf{Build.}
\begin{enumerate}
\item Compute the frequency table $F$ of the first $k$ letters of $A$ in $O(k)$ time.
\item Slide the window one step at a time: to move from the window starting at $A[i]$ to the
one starting at $A[i+1]$, subtract one from the count of letter $A[i]$ and add one to the count of
letter $A[i+k]$. This is $O(1)$ per step.
\item After each step, look up the current table in a hash table $H$ and add one to its count
(or insert it with count 1). Each table has 26 counts, each at most $|A|$, so it is a
constant number of words and can be hashed and compared in $O(1)$.
\end{enumerate}
There are $|A|-k+1$ windows, so building takes expected $O(k+|A|)=O(|A|)$ time.

\textbf{Query($B$), $|B|=k$.} Compute the frequency table $f$ of $B$ in $O(k)$ time and look it up
in $H$ in expected $O(1)$ time. Return the stored count, or 0 if $f$ is not in $H$.

\textbf{Correctness.} A substring of $A$ of length $k$ is an anagram of $B$ if and only if it has
the same frequency table as $B$. $H[f]$ counts exactly those windows. The sliding update gives the
correct table for each window because it changes only the two letters that enter and leave.

\textbf{Running time.} Build is expected $O(|A|)$ and a query is expected $O(k)$.

\problempart % Problem 5b
Build the data structure from part (a) on $T$ with window length $k$ (all $S_i$ have length $k$), in
expected $O(|T|)$ time. Then query each $S_i$ and store the answers $a_i$ in an array: $n$
queries of expected $O(k)$ each, so expected $O(nk)$. Total expected $O(|T|+nk)$. It is correct because each
query returns the anagram substring count of $S_i$ in $T$, by part (a).

\problempart Submit your implementation to {\small\url{alg.mit.edu}}.
\end{problemparts}

\end{problems}

\end{document}
